\section{All Rankin cross products among the known comparison constituents}
\label{sec:s5-rankin-cross}

The following calculation tests every pair among the specified
five-dimensional comparison and its sign twist and the known
four- and six-dimensional cuspidal factors. It uses the
Rankin--Selberg functions of those actual cuspidal representations;
it assumes no automorphic tensor-product lift.

Retain the actual point \(\gamma\), the positive integer \(k\),
and the nonnegative orders of (MA-RB.46). Set
\[
 \begin{gathered}
 (\pi_1,\pi_2,\pi_3,\pi_4,\pi_5)
   =(\varpi,\varpi\otimes\eta,\pi_V,\pi_{V\eta},\pi_X),\\
 (U_1,U_2,U_3,U_4,U_5)=(W,W\eta,V,V\eta,X),\\
 \Theta=b\eta+v_0V+v_1V\eta+(n+k)W+x_1W\eta+eX,\qquad
 q=v_1+e-k\ge0.
 \end{gathered}
 \tag{MA-RB.61}
\]
Here \(\pi_V,\pi_{V\eta},\pi_X\) are the established
automorphic realizations of those Artin representations.
In particular \(\pi_V\) is distinct notation from the
four-dimensional symplectic comparison \(\Pi\) in MA-RB-005.
All five \(\pi_i\) are unitary cuspidal and self-dual.
For \(\varpi\) this was proved in MA-S5-002; sign twisting
preserves it. For the other three it follows from their
actual Artin correspondence and the real self-dual
representations, or from their local parameters and strong
multiplicity one.

\paragraph{MA-RB-006 (the complete fifteen-pair order constraints).}
Choose \(S\) to contain infinity and all ramified or omitted
places in this finite family, and retain \(D^S\) from
MA-RB-005. Put
\[
 \mathcal L_{ij}(s)=L^S(s,\pi_i\times\pi_j^\vee),\qquad
 E=\begin{pmatrix}
 2&-2&0&0&0\\
 -2&2&0&0&0\\
 0&0&0&0&0\\
 0&0&0&0&0\\
 0&0&0&0&0
 \end{pmatrix}.
 \tag{MA-RB.62}
\]
The exact meromorphic identities are
\[
          \mathcal L_{ij}(s)
             =(D^S(s))^{E_{ij}}L^S(s,U_i\otimes U_j)
                     \qquad(1\le i,j\le5).
 \tag{MA-RB.63}
\]
Each \(\mathcal L_{ij}\) is holomorphic at \(\gamma\), and
its order is
\[
       r_{ij}
          =\langle U_i\otimes U_j,\Theta\rangle_{S_5}
                       -E_{ij}k .
 \tag{MA-RB.64}
\]
Let
\[
 M=\begin{pmatrix}
 0&4&1&2&2\\
 4&0&2&1&2\\
 1&2&1&1&2\\
 2&1&1&1&2\\
 2&2&2&2&3
 \end{pmatrix}.
 \tag{MA-RB.65}
\]
Then \(r_{ij}\ge kM_{ij}\) entry by entry. Every one of
the fifteen holomorphy inequalities \(r_{ij}\ge0\) is
already implied by the nonnegative orders in (MA-RB.61)
and the earlier inequality \(q\ge0\). The formal
assignment \(\Theta=k(V\eta+W)\), \(n=0\), attains
\((r_{ij})=kM\) and regular order \(9k\).
Thus this complete family of Rankin holomorphy constraints
neither excludes that formal assignment nor improves the
previous normal-closure multiplicity bound.

\begin{proof}
\emph{The local comparison for all fifteen pairs.}
At an exceptional prime the class is \(4A\).
The five-dimensional parameter multisets \(T_A,T_W\)
are (MA-RB.55). For their \(j\)th power traces \(a_j,w_j\),
\[
 a_j-w_j=
  \begin{cases}4,&j\text{ odd},\\0,&j\text{ even},\end{cases}
 \quad
 a_j^2-w_j^2=
  \begin{cases}8,&j\text{ odd},\\0,&j\text{ even}.\end{cases}
 \tag{MA-RB.66}
\]
Indeed \(a_j=3,w_j=-1\) for odd \(j\), while for even
\(j\) both are \(3+i^j+(-i)^j\).
The sign character equals \(-1\) here, so the traces
for \(\varpi\otimes\eta,W\eta\) are
\((-1)^j a_j,(-1)^jw_j\).

The local standard parameter for \(V\) at \(4A\) is
\(\{1,-1,i,-i\}\). Its sign twist has the same multiset,
and the parameter of \(X=\wedge^2V\) is
\(\{1,-1,i,i,-i,-i\}\). Consequently the \(j\)th
power trace of each of \(V,V\eta,X\) is zero whenever
\(j\) is odd. This can also be read directly from the
character table: every odd power of a \(4A\) element
remains in \(4A\).

For local multisets \(P,Q\), the logarithm of the
Rankin Euler factor is
\(\sum_{j\ge1}\operatorname{tr}(P^j)
\operatorname{tr}(Q^j)T^j/j\).
All our local multisets are self-dual, so the use of
\(\pi_j^\vee\) changes none of these calculations.
In a mixed comparison with \(V,V\eta\), or \(X\),
the trace difference vanishes for even \(j\) by
(MA-RB.66), and for odd \(j\) by the zero trace of
the second factor. Its complete Euler ratio is
therefore one. The known three representations
already match at all unramified primes, so a pair
entirely among them has ratio one as well.

For the two diagonal five-dimensional pairs the
logarithmic ratio is
\[
       8\sum_{\substack{j\ge1\\j\ {\rm odd}}}\frac{T^j}{j}
                     =2\log\Delta(T).
 \tag{MA-RB.67}
\]
For the cross pair \((1,2)\), its extra factor
\((-1)^j\) negates this expression, giving
\(\Delta(T)^{-2}\). Exponentiating the absolutely
convergent logarithms near \(T=0\) proves the exact
rational local identities. Outside \(\Omega\) all
five original unramified parameters already match
their corresponding Artin ones.
Multiplying the local identities proves (MA-RB.63)
in an absolute-convergence half-plane, and hence
meromorphically. This proof determines the complete
local polynomials, not only their first coefficients.

\emph{The complete finite tensor calculation.}
For clarity the following table gives the multiplicities
of every irreducible in \(U_i\otimes U_j\); symmetric
entries have been listed once.
\[
 \begin{array}{c|rrrrrrr}
  ij&\mathbf1&\eta&V&V\eta&W&W\eta&X\\ \hline
  11&1&0&1&1&1&1&1\\
  12&0&1&1&1&1&1&1\\
  13&0&0&1&0&1&1&1\\
  14&0&0&0&1&1&1&1\\
  15&0&0&1&1&1&1&2\\
  22&1&0&1&1&1&1&1\\
  23&0&0&0&1&1&1&1\\
  24&0&0&1&0&1&1&1\\
  25&0&0&1&1&1&1&2\\
  33&1&0&1&0&1&0&1\\
  34&0&1&0&1&0&1&1\\
  35&0&0&1&1&1&1&1\\
  44&1&0&1&0&1&0&1\\
  45&0&0&1&1&1&1&1\\
  55&1&1&1&1&2&2&1
 \end{array}
 \tag{MA-RB.68}
\]
Here is an exact method proving every entry, rather
than an assumed tensor-product rule. In class order
\(1,2A,2B,3A,6A,4A,5A\), take the already checked
character-table columns and class sizes
\((1,10,15,20,20,30,24)\). Multiply the two selected
columns entry by entry. Its coefficient in column
\(\chi\) is
\[
 \frac1{120}\sum_{\mathcal C}|\mathcal C|\,
      \chi_{U_i}(\mathcal C)\chi_{U_j}(\mathcal C)
                                      \chi(\mathcal C).
\]
Substitution gives (MA-RB.68), and multiplying the
resulting integer rows back by the character table
recovers each product column. Equality of characters
then proves every displayed decomposition.
The old \(W\otimes W\) identity is included for the
complete matrix, using the same previously checked table.
All remaining entries also follow from the displayed
products and \(X\eta=X\), but no sign-twist convention
has been left implicit in (MA-RB.68).

\emph{Actual holomorphy and pointwise orders.}
The Rankin--Selberg theorem for two unitary cuspidal
representations is applicable to every \(\pi_i,\pi_j\).
For unequal dimensions the completed function is entire.
For equal dimensions its possible poles lie on
\(\operatorname{Re}s=0\) or \(1\), and thus away
from the present \(\gamma\). Deleting finite Euler
and archimedean factors introduces no pole: their
reciprocals are entire. In this finite-image situation
they are in fact units at a real point of \((0,1)\),
by the same local-parameter argument as MA-S5-003.
Hence every \(\mathcal L_{ij}\) is holomorphic at
\(\gamma\), as asserted.

On the Artin side all deleted finite factors are
units there. The order of the trivial constituent
is zero because \(\zeta_{\mathbb Q}(\gamma)\ne0\).
Artin multiplicativity and (MA-RB.61) therefore give
the inner product in (MA-RB.64). Also
\(\operatorname{ord}_\gamma D^S=-k\). Taking orders
in the genuine meromorphic identity (MA-RB.63)
proves (MA-RB.64) exactly, including its sign.
No general holomorphy of the unknown \(L_W\) was
assumed: its positive order \(n+k\) at this particular
point is the consequence already proved in MA-RB-002.

\emph{All inequalities and the formal witness.}
The two nonzero positive exponents in \(E\) give
\[
 r_{11}=r_{22}
     =v_0+v_1+n+k+x_1+e-2k
     =q+n+v_0+x_1\ge0 .
 \tag{MA-RB.69}
\]
The only negative exponents are \(E_{12}=E_{21}=-2\);
they add \(2k\) to an already nonnegative genuine-character
pairing. All other exponents vanish, so their
nonnegativity is immediate from the genuine tensor
characters and the nonnegative orders in \(\Theta\).
Thus every holomorphy inequality follows from the
prior order constraints.

More explicitly, subtracting \(kM_{ij}\) from
(MA-RB.64) gives the following exhaustive table:
\[
 \begin{array}{c|l}
 ij&r_{ij}-kM_{ij}\\ \hline
 11,22&q+n+v_0+x_1\\
 12&b+q+n+v_0+x_1\\
 13,24&e+n+v_0+x_1\\
 14,23&q+n+x_1\\
 15,25&e+q+n+v_0+x_1\\
 33,44&e+n+v_0\\
 34&b+q+x_1\\
 35,45&q+n+v_0+x_1\\
 55&b+q+2n+v_0+2x_1 .
 \end{array}
 \tag{MA-RB.70}
\]
Every entry is nonnegative, proving the stated
entrywise matrix bound. Under the formal assignment
\[
        (b,n,v_0,v_1,e,x_1)=(0,0,0,k,0,0)
\]
one has \(q=0\), so every entry in (MA-RB.70)
vanishes. Hence \(r=kM\), while the regular-character
order remains \(9k\), as in MA-RB-004.
The same tuple also satisfies all four quadratic
Schur inequalities of MA-RB-005. It follows that
adding all fifteen Rankin order nonnegativities
cannot by itself shrink that formal feasible set
or strengthen the bound it attains.
\end{proof}

The conclusion concerns exactly the inequalities supplied
by these known Rankin holomorphy theorems. It does not
assert that arbitrary formal order matrices are
arithmetically realizable, or rule out further analytic
restrictions on their simultaneous zeros. In particular
the numerous actual zeros forced by (MA-RB.65) could be
excluded by a suitable additional nonvanishing theorem;
the holomorphy theorem used here supplies none.
The residual shared-real-zero criterion of MA-RB-002
therefore remains open.

\paragraph{Primary inputs and verification scope.}
The unramified exceptional parameters and the real
\(S_5\) table are the same Calegari \S2 and
Theorem~4.7 inputs used in MA-RB-005. Its four- and
six-dimensional automorphy is Theorem~1.1, already
checked in MA-S5. The argument reuses those
correspondences and does not assume the missing
five-dimensional comparison.

For the analytic input, Cogdell (2004), Lecture~9,
Theorems~9.1(i) and 9.2(i), printed p.~74
(PDF p.~78), state the precise unequal- and
equal-dimension conclusions. The complete terminal
continuation and pole argument in Lecture~9,
\S\S2--3, printed p.~71 (PDF p.~75), was freshly
read for this application. It expresses the global
function through the global integrals and entire
local ratios, then chooses the local ratios nonzero
at the point in question. The underlying global
integral and local Whittaker theorems remain declared
classical inputs; their full earlier proofs are not
newly audited here.

The complete exact calculation is reproducible in
\texttt{checks/s5/s5\_rankin\_cross.py} and its JSON
output. It checks all fifteen product-character
decompositions, the full cyclotomic local ratios,
and the general-\(k\) order matrix. It reuses the
existing checked table and makes no arithmetic
zero-location computation.

Cogdell, J. W. (2004). Lectures on \(L\)-functions,
converse theorems, and functoriality for
\(\mathrm{GL}_n\). In \emph{Lectures on automorphic
\(L\)-functions} (Fields Institute Monographs,
Vol.~20). American Mathematical Society.
\url{https://people.math.osu.edu/cogdell.1/fields-www.pdf}.
